\section{Planteamiento del Problema}

En la Facultad de Ciencias y Tecnología de la Universidad Mayor de San Simón, la formación en seguridad informática se aborda principalmente en niveles avanzados mediante la asignatura de Seguridad de Sistemas. Esta materia es una instancia importante para desarrollar competencias relacionadas con la identificación, análisis y respuesta ante amenazas digitales.

Sin embargo, el desarrollo de la asignatura mantiene un enfoque predominantemente teórico. Las prácticas suelen realizarse de forma autónoma por los estudiantes, fuera del aula o mediante ejercicios complementarios, lo que reduce las posibilidades de experimentación guiada, seguimiento y retroalimentación durante el proceso de aprendizaje.

Además, la cantidad de estudiantes inscritos dificulta el acompañamiento individual. En la gestión 2026-I se registran aproximadamente 130 estudiantes, mientras que en la gestión 2025-II se alcanzaron alrededor de 180. Esta situación limita la supervisión de prácticas y refuerza evaluaciones centradas en exámenes teóricos y trabajos expositivos.

La ciberseguridad requiere que el estudiante no solo comprenda conceptos, sino que practique la toma de decisiones frente a escenarios de riesgo. Cuando existe poca práctica guiada, se genera una brecha entre el conocimiento conceptual y la capacidad de aplicarlo ante situaciones similares a las de un entorno profesional.

Frente a esta situación, una herramienta basada en simulación gamificada puede complementar la enseñanza teórica mediante escenarios controlados de amenazas digitales, retos progresivos y retroalimentación sobre las decisiones del estudiante. De esta manera, se busca fortalecer habilidades prácticas sin comprometer sistemas reales.

\subsection*{Formulación del problema}
\addcontentsline{toc}{subsection}{Formulación del problema}

\noindent ¿De qué manera el desarrollo de un videojuego de simulación gamificada, basado en escenarios de amenazas digitales y retroalimentación por escenario, contribuye al fortalecimiento de las habilidades prácticas de respuesta ante incidentes en estudiantes de Ingeniería Informática y Sistemas de la UMSS?
